\documentclass[10pt,a4paper]{amsart}
\usepackage[margin=19mm]{geometry}
\usepackage{amsmath,amssymb,mathtools}
\usepackage[scr=rsfs,cal=euler]{mathalfa}
\usepackage{xfrac}
\usepackage[unicode,hidelinks,bookmarks=false]{hyperref}
\newcommand{\ie}{i.e.\ }
\newcommand{\cf}{cf.\ }
\newcommand{\resp}{respectively\ }
\newcommand{\Subsection}{\subsection}
\providecommand{\nobreakdash}{\nobreak\hbox{-}}
\setlength{\emergencystretch}{2em}
\multlinegap0pt
\usepackage{amssymb}
\usepackage{mathtools}
\usepackage{enumerate}
\usepackage{dsfont}
\usepackage{float}
\usepackage{tikz}
\newtheorem{lemma}{Lemma}
\title{Spectral analysis of multivariate stationary Hawkes processes}
\author{Yifu Tang, Conor Kresin, Boris Baeumer, Ting Wang}
\date{Source: arXiv:2604.10376v1 (CC BY 4.0)}
\begin{document}
\maketitle

\noindent\emph{Source and licence.} Spectral analysis of multivariate stationary Hawkes processes. Version arXiv:2604.10376v1 (CC BY 4.0), distributed by arXiv under the Creative Commons Attribution 4.0 International licence, http://creativecommons.org/licenses/by/4.0/, as recorded in the arXiv licence metadata for this exact version and shown on the abstract page https://arxiv.org/abs/2604.10376v1. Full licence notice: \texttt{licenses/arxiv-2604.10376-CC-BY-4.0.txt}.\par\smallskip

\noindent\textit{Editorial scope.} Counted unit: the article's lemma lem:indec-connected (source0.tex lines 1926--1928). The article's complete original proof of it is delivered with it (source0.tex lines 1930--1939, one \texttt{proof} environment of 10 lines). No mathematics is written, repaired or re-derived: the statement and the proof are the article's own lines, in the article's own notation, and no retained line is substituted or re-flowed.  No further source-local context is needed: every cross-reference of every retained line resolves inside the two delivered spans.  Nothing was omitted from any retained span, so the package carries no omission record.  No retained line contains a citation, so the wrapper carries no bibliography and no bibliography entry.  No retained line contains a figure environment, an image-inclusion command or drawing code, so no image is delivered and the package declares no asset dependency.  Editorial formatting only, in addition: an AMS \texttt{amsart} preamble that restates the article's own declarations and macros used by the retained lines, namely the environments \texttt{lemma} on separate counters, together with the standard packages those lines need.  The retained environments print in this document as Lemma 1 in order of appearance, whereas the article prints the same items with the numbering of its own full document.  The complete native source of the article as distributed in the arXiv e-print for this exact version is bundled unchanged as \texttt{sources/198269/source0.tex}.
\begin{lemma}[Indecomposable $\iff$ connected]\label{lem:indec-connected}
Let $G$ be the graph defined above (ignore loops when talking about connectivity). Then the partition is indecomposable if and only if $G$ is connected.
\end{lemma}

\begin{proof}
\emph{($\Rightarrow$)} Suppose, towards a contradiction, that $G$ is disconnected. Take the vertex set $V_1$ of any nonempty proper connected component. Because edges never cross between components, every row $R_i=\{p_i,-p_i\}$ is either entirely supported on blocks indexed by $V_1$ or entirely on blocks indexed by the complement $V_1^c$. Therefore
\[
\bigcup_{j\in V_1} Q_j \;=\; \bigcup_{i: R_i\subset \cup_{j\in V_1}Q_j} R_i
\]
is a union of whole rows, contradicting indecomposability.\\


\emph{($\Leftarrow$)} Conversely, suppose some strict subfamily $\{Q_{j_1},\dots,Q_{j_s}\}$ equals a union of whole rows. Then no row $R_i$ is split across this subfamily and its complement; equivalently, there is no edge in $G$ joining $\{j_1,\dots,j_s\}$ to its complement. Hence $G$ is disconnected.
\end{proof}

\end{document}
